\chapter{1996年全国卷（文）}
\section{单选题}

\begin{problem}
设全集 \(I = \{1, 2, 3, 4, 5, 6, 7\}\)，集合 \(A = \{1, 3, 5, 7\}\)，\(B = \{3, 5\}\)，则（\quad）

\choices
    {\(I = A\cup B\)}
    {\(I = \overline{A}\cup B\)}
    {\(I = A\cup \overline{B}\)}
    {\(I = \overline{A}\cup \overline{B}\)}
\end{problem}

\begin{answer}
C
\end{answer}

\begin{solution}
由全集 \(I\) 可得 \(\overline{A}=\{2,4,6\}\)，\(\overline{B}=\{1,2,4,6,7\}\). 因此
\[
A\cup\overline{B}=\{1,3,5,7\}\cup\{1,2,4,6,7\}=I
\]
而 \(A\cup B=A\ne I\)，\(\overline{A}\cup B=\{2,3,4,5,6\}\ne I\)，\(\overline{A}\cup\overline{B}=\{1,2,4,6,7\}\ne I\)，故选 C.
\end{solution}

\begin{problem}
当 \(a > 1\) 时，在同一坐标系中，函数 \(y = a^{-x}\) 与 \(y = \log_a x\) 的图象（\quad）

\choices
    {\choicebitmap[width=0.15\paperwidth]{img/1996/national_liberal/q2_optA.png}}
    {\choicebitmap[width=0.15\paperwidth]{img/1996/national_liberal/q2_optB.png}}
    {\choicebitmap[width=0.15\paperwidth]{img/1996/national_liberal/q2_optC.png}}
    {\choicebitmap[width=0.15\paperwidth]{img/1996/national_liberal/q2_optD.png}}
\end{problem}

\begin{answer}
A
\end{answer}

\begin{solution}
因为 \(a>1\)，所以 \(y=a^{-x}\) 的定义域为 \(\R\)，在 \(\R\) 上单调递减，并且经过点 \((0,1)\). 函数 \(y=\log_a x\) 的定义域为 \((0,+\infty)\)，在其定义域内单调递增，并且经过点 \((1,0)\)，当 \(x\to0^+\) 时函数值趋于 \(-\infty\). 两条曲线在第一象限内相交，且一条递减、一条递增，图形与选项 A 相符，故选 A.
\end{solution}

\begin{problem}
若 \(\sin^2 x > \cos^2 x\)，则 \(x\) 的取值范围是（\quad）

\choices
    {\(\left\{x\mid 2k\pi -\frac{3}{4}\pi < x < 2k\pi +\frac{1}{4}\pi ,k\in \Z\right\}\)}
    {\(\left\{x\mid 2k\pi +\frac{1}{4}\pi < x < 2k\pi +\frac{5}{4}\pi ,k\in \Z\right\}\)}
    {\(\left\{x\mid k\pi -\frac{1}{4}\pi < x < k\pi + \frac{1}{4}\pi ,k\in \Z\right\}\)}
    {\(\left\{x\mid k\pi +\frac{1}{4}\pi < x < k\pi +\frac{3}{4}\pi ,k\in \Z\right\}\)}
\end{problem}

\begin{answer}
D
\end{answer}

\begin{solution}
由 \(\sin^2x-\cos^2x=-\cos2x\)，题设等价于 \(\cos2x<0\). 因此存在 \(k\in\Z\)，使
\[
\frac{\pi}{2}+2k\pi<2x<\frac{3\pi}{2}+2k\pi
\]
两边同除以 \(2\)，得
\(k\pi+\frac{\pi}{4}<x<k\pi+\frac{3\pi}{4}\)，\(k\in\Z\)
故选 D.
\end{solution}

\begin{problem}
复数 \(\frac{\left(2 + 2\i\right)^4}{\left(1 - \sqrt{3}\i\right)^5}\) 等于（\quad）

\choices
    {\(1 + \sqrt{3}\i\)}
    {\(-1 + \sqrt{3}\i\)}
    {\(1 - \sqrt{3}\i\)}
    {\(-1 - \sqrt{3}\i\)}
\end{problem}

\begin{answer}
B
\end{answer}

\begin{solution}
写成三角形式，有 \(2+2\i=2\sqrt{2}\left(\cos\frac{\pi}{4}+\i\sin\frac{\pi}{4}\right)\)，所以
\[
(2+2\i)^4=\left(2\sqrt{2}\right)^4\left(\cos\pi+\i\sin\pi\right)=-64
\]
又 \(1-\sqrt{3}\i=2\left(\cos\left(-\frac{\pi}{3}\right)+\i\sin\left(-\frac{\pi}{3}\right)\right)\)，从而
\[
(1-\sqrt{3}\i)^5=32\left(\cos\frac{\pi}{3}+\i\sin\frac{\pi}{3}\right)=16+16\sqrt{3}\i
\]
于是
\[
\frac{(2+2\i)^4}{(1-\sqrt{3}\i)^5}=\frac{-64}{16+16\sqrt{3}\i}=-\frac{4}{1+\sqrt{3}\i}=-\frac{4(1-\sqrt{3}\i)}{4}=-1+\sqrt{3}\i
\]
故选 B.
\end{solution}

\begin{problem}
\(6\text{ 名}\)同学排成一排，其中甲，乙两人必须排在一起的不同排法有（\quad）

\choices
    {\(720\text{ 种}\)}
    {\(360\text{ 种}\)}
    {\(240\text{ 种}\)}
    {\(120\text{ 种}\)}
\end{problem}

\begin{answer}
C
\end{answer}

\begin{solution}
把甲、乙两人看作一个整体，与另外四人一起排列，共有 \(5!\) 种排法. 甲、乙在这个整体中的顺序有 \(2!\) 种，因此不同排法共有
\[
5!\times2!=120\times2=240
\]
种，故选 C.
\end{solution}

\begin{problem}
已知 \(\alpha\) 是第三象限角且 \(\sin \alpha = -\frac{24}{25}\)，则 \(\tan \frac{\alpha}{2} =\)（\quad）

\choices
    {\(\frac{4}{3}\)}
    {\(\frac{3}{4}\)}
    {\(-\frac{3}{4}\)}
    {\(-\frac{4}{3}\)}
\end{problem}

\begin{answer}
D
\end{answer}

\begin{solution}
因为 \(\alpha\) 是第三象限角，所以 \(\cos\alpha<0\). 由 \(\sin^2\alpha+\cos^2\alpha=1\) 得
\[
\cos\alpha=-\sqrt{1-\left(-\frac{24}{25}\right)^2}=-\frac{7}{25}
\]
利用半角公式
\[
\tan\frac{\alpha}{2}=\frac{\sin\alpha}{1+\cos\alpha}
=\frac{-\frac{24}{25}}{1-\frac{7}{25}}
=\frac{-\frac{24}{25}}{\frac{18}{25}}=-\frac{4}{3}
\]
故选 D.
\end{solution}

\begin{problem}
如果直线 \(l\)，\(m\) 与平面 \(\alpha\)，\(\beta\)，\(\gamma\) 满足：\(l = \beta \cap \gamma\)，\(l \parallel \alpha\)，\(m \subset \alpha\)，\(m \perp \gamma\)，那么必有（\quad）

\choices
    {\(\alpha \perp \gamma\) 且 \(l \perp m\)}
    {\(\alpha \perp \gamma\) 且 \(m \parallel \beta\)}
    {\(m \parallel \beta\) 且 \(l \perp m\)}
    {\(\alpha \parallel \beta\) 且 \(\alpha \perp \gamma\)}
\end{problem}

\begin{answer}
A
\end{answer}

\begin{solution}
由 \(m\perp\gamma\) 且 \(m\subset\alpha\)，可知平面 \(\alpha\perp\gamma\). 又因为 \(l\subset\gamma\) 且 \(l\parallel\alpha\)，所以 \(l\) 的方向向量既平行于平面 \(\alpha\)，又在平面 \(\gamma\) 内，因而平行于两平面的交线方向. 直线 \(m\perp\gamma\)，所以 \(l\) 的方向向量与 \(m\) 的方向向量垂直；按两直线所成角的定义，得 \(l\perp m\). 因此必有 \(\alpha\perp\gamma\) 且 \(l\perp m\)，故选 A.
\end{solution}

\begin{problem}
当 \(-\frac{\pi}{2} \leqslant x \leqslant \frac{\pi}{2}\)，\(f(x) = \sin x + \sqrt{3} \cos x\) 的（\quad）

\choices
    {最大值是 \(1\)，最小值是 \(-1\)}
    {最大值是 \(1\)，最小值是 \(-\frac{1}{2} \)}
    {最大值是 \(2\)，最小值是 \(-2\)}
    {最大值是 \(2\)，最小值是 \(-1\)}
\end{problem}

\begin{answer}
D
\end{answer}

\begin{solution}
利用和角公式，得
\[
f(x)=\sin x+\sqrt{3}\cos x=2\sin\left(x+\frac{\pi}{3}\right)
\]
当 \(-\frac{\pi}{2}\leqslant x\leqslant\frac{\pi}{2}\) 时，\(-\frac{\pi}{6}\leqslant x+\frac{\pi}{3}\leqslant\frac{5\pi}{6}\). 在此区间内，正弦函数的最大值为 \(1\)，在 \(x+\frac{\pi}{3}=\frac{\pi}{2}\) 时取得；最小值为 \(-\frac{1}{2}\)，在左端点取得. 因而 \(f(x)\) 的最大值为 \(2\)，最小值为 \(-1\)，故选 D.
\end{solution}

\begin{problem}
中心在原点，准线方程为 \(x = \pm 4\)，离心率为 \(\frac{1}{2}\) 的椭圆方程是（\quad）

\choices
    {\(\frac{x^2}{4} +\frac{y^2}{3} = 1\)}
    {\(\frac{x^2}{3} +\frac{y^2}{4} = 1\)}
    {\(\frac{x^2}{4} + y^2 = 1\)}
    {\(x^2 + \frac{y^2}{4} = 1\)}
\end{problem}

\begin{answer}
A
\end{answer}

\begin{solution}
设椭圆的长半轴为 \(a_0\)，短半轴为 \(b_0\)，离心距为 \(c_0\). 由于准线为 \(x=\pm4\)，长轴在 \(x\) 轴上，且椭圆准线满足 \(\frac{a_0}{e}=4\). 已知 \(e=\frac12\)，所以 \(a_0=2\). 再由 \(e=\frac{c_0}{a_0}\) 及 \(b_0^2=a_0^2-c_0^2=a_0^2(1-e^2)\)，得 \(b_0^2=4\left(1-\frac14\right)=3\). 因此椭圆方程为
\[
\frac{x^2}{4}+\frac{y^2}{3}=1
\]
故选 A.
\end{solution}

\begin{problem}
圆锥母线长为 \(1\)，侧面展开图圆心角为 \(240^{\circ}\)，该圆锥的体积（\quad）

\choices
    {\(\frac{2\sqrt{2}\pi}{81}\)}
    {\(\frac{8\pi}{81}\)}
    {\(\frac{4\sqrt{5}\pi}{81}\)}
    {\(\frac{10\pi}{81}\)}
\end{problem}

\begin{answer}
C
\end{answer}

\begin{solution}
圆锥侧面展开图是半径为母线长 \(1\) 的扇形，其弧长等于圆锥底面周长. 因此底面半径 \(r\) 满足
\[
2\pi r=\frac{240^{\circ}}{360^{\circ}}\cdot2\pi\cdot1=\frac{4\pi}{3}
\]
解得 \(r=\frac{2}{3}\). 圆锥高为
\[
h=\sqrt{1^2-r^2}=\sqrt{1-\frac{4}{9}}=\frac{\sqrt{5}}{3}
\]
所以体积
\[
V=\frac{1}{3}\pi r^2h=\frac{1}{3}\pi\left(\frac{2}{3}\right)^2\frac{\sqrt{5}}{3}=\frac{4\sqrt{5}\pi}{81}
\]
故选 C.
\end{solution}

\begin{problem}
椭圆 \(25x^{2} - 150x + 9y^{2} + 18y + 9 = 0\) 的两个焦点坐标是（\quad）

\choices
    {\(\left(-3, 5\right)\)，\(\left(-3, -3\right)\)}
    {\(\left(3, 3\right)\)，\(\left(3, -5\right)\)}
    {\(\left(1, 1\right)\)，\(\left(-7, 1\right)\)}
    {\(\left(7, -1\right)\)，\(\left(-1, -1\right)\)}
\end{problem}

\begin{answer}
B
\end{answer}

\begin{solution}
配方得
\[
25(x-3)^2+9(y+1)^2=225
\]
即
\[
\frac{(x-3)^2}{9}+\frac{(y+1)^2}{25}=1
\]
椭圆中心为 \((3,-1)\)，长半轴为 \(5\)，短半轴为 \(3\)，故半焦距为 \(\sqrt{5^2-3^2}=4\). 由于长轴沿 \(y\) 轴方向，两个焦点为
\((3,-1+4)=(3,3)\)，\((3,-1-4)=(3,-5)\)
故选 B.
\end{solution}

\begin{problem}
将边长为 \(a\) 的正方形 \(ABCD\) 沿对角线 \(AC\) 折起，使得 \(BD = a\)，则三棱锥 \(DABC\) 的体积为（\quad）

\choices
    {\(\frac{a^3}{6}\)}
    {\(\frac{a^3}{12}\)}
    {\(\frac{\sqrt{3}}{12} a^3\)}
    {\(\frac{\sqrt{2}}{12} a^3\)}
\end{problem}

\begin{answer}
D
\end{answer}

\begin{solution}
设折起后两平面的二面角为 \(\theta\)，令 \(M\) 为对角线 \(AC\) 的中点. 在正方形中，\(BM\perp AC\)，\(DM\perp AC\)，且
\[
BM=DM=\frac{a}{\sqrt{2}}
\]
折起前 \(B\)，\(D\) 在直线 \(AC\) 的两侧，所以折起后 \(\angle BMD=\pi-\theta\). 因而
\[
BD^2=BM^2+DM^2-2BM\cdot DM\cos(\pi-\theta)=a^2(1+\cos\theta)
\]
由 \(BD=a\) 得 \(\cos\theta=0\)，所以 \(\theta=\frac{\pi}{2}\). 点 \(D\) 到平面 \(ABC\) 的距离为
\[
DM\sin\theta=\frac{a}{\sqrt{2}}
\]
以 \(\triangle ABC\) 为底面，其面积为 \(\frac{a^2}{2}\)，故三棱锥体积为
\[
V=\frac{1}{3}\cdot\frac{a^2}{2}\cdot\frac{a}{\sqrt{2}}=\frac{\sqrt{2}}{12}a^3
\]
故选 D.
\end{solution}

\begin{problem}
等差数列 \(\{a_{n}\}\) 的前 \(m\) 项和为 \(30\)，前 \(2m\) 项和为 \(100\)，则它的前 \(3m\) 项和为（\quad）

\choices
    {\(130\)}
    {\(170\)}
    {\(210\)}
    {\(260\)}
\end{problem}

\begin{answer}
C
\end{answer}

\begin{solution}
把等差数列分成连续的三个长度均为 \(m\) 的数段，记三段的和分别为 \(T_1\)，\(T_2\)，\(T_3\). 已知 \(T_1=S_m=30\)，且
\[
T_1+T_2=S_{2m}=100
\]
所以 \(T_2=70\). 第二段的每一项比第一段对应项多 \(md\)，其中 \(d\) 是公差，因此
\[
T_2-T_1=m\cdot md=m^2d
\]
第三段与第二段的对应项之差同样是 \(md\)，从而 \(T_3-T_2=m^2d=T_2-T_1=40\)，得 \(T_3=110\). 因此
\[
S_{3m}=T_1+T_2+T_3=30+70+110=210
\]
故选 C.
\end{solution}

\begin{problem}
设双曲线 \(\frac{x^2}{a^2} - \frac{y^2}{b^2} = 1\)（\(0 < a < b\)）的半焦距为 \(c\)，直线 \(l\) 过 \(\left(a, 0\right)\)，\(\left(0, b\right)\) 两点，已知原点到直线 \(l\) 的距离为 \(\frac{\sqrt{3}}{4} c\)，则双曲线的离心率为（\quad）

\choices
    {\(2\)}
    {\(\sqrt{3}\)}
    {\(\sqrt{2}\)}
    {\(\frac{2\sqrt{3}}{3}\)}
\end{problem}

\begin{answer}
A
\end{answer}

\begin{solution}
直线 \(l\) 过 \((a,0)\) 和 \((0,b)\)，故其方程为
\[
\frac{x}{a}+\frac{y}{b}=1
\]
即 \(bx+ay-ab=0\). 原点到此直线的距离为
\[
\frac{ab}{\sqrt{a^2+b^2}}
\]
双曲线的半焦距满足 \(c^2=a^2+b^2\)，所以由题设
\[
\frac{ab}{c}=\frac{\sqrt{3}}{4}c
\]
即 \(4ab=\sqrt{3}(a^2+b^2)\). 令 \(t=\frac{b}{a}>1\)，得
\[
\sqrt{3}t^2-4t+\sqrt{3}=0
\]
解得 \(t=\sqrt{3}\) 或 \(t=\frac{1}{\sqrt{3}}\)，由 \(t>1\) 取 \(t=\sqrt{3}\). 因此双曲线的离心率为
\[
e=\frac{c}{a}=\sqrt{1+\frac{b^2}{a^2}}=\sqrt{1+t^2}=2
\]
故选 A.
\end{solution}

\begin{problem}
设 \(f\left(x\right)\) 是 \(\left(-\infty, +\infty\right)\) 上的奇函数，\(f\left(x + 2\right) = -f\left(x\right)\)，当 \(0 \leqslant x \leqslant 1\) 时，\(f\left(x\right) = x\)，则 \(f\left(7.5\right)\) 等于（\quad）

\choices
    {\(0.5\)}
    {\(-0.5\)}
    {\(1.5\)}
    {\(-1.5\)}
\end{problem}

\begin{answer}
B
\end{answer}

\begin{solution}
由 \(f(x+2)=-f(x)\)，再把 \(x\) 换成 \(x+2\)，得 \(f(x+4)=f(x)\)，所以 \(f\) 的周期为 \(4\). 因此
\[
f(7.5)=f(-0.5)
\]
又因为 \(f\) 是奇函数，\(f(-0.5)=-f(0.5)\)，而 \(0\leqslant0.5\leqslant1\)，所以 \(f(0.5)=0.5\). 故
\[
f(7.5)=-0.5
\]
故选 B.
\end{solution}

\section{填空题}

\begin{problem}
已知点 \((-2,3)\) 与抛物线 \(y^{2} = 2px\) （\(p > 0\)）的焦点的距离是 \(5\)，则 \(p =\)\fillinblank{}.
\end{problem}

\begin{answer}
\(4\)
\end{answer}

\begin{solution}
抛物线 \(y^2=2px\) 的焦点为 \(F\left(\frac{p}{2},0\right)\). 由点 \((-2,3)\) 到焦点的距离为 \(5\)，有
\[
\left(-2-\frac{p}{2}\right)^2+3^2=5^2
\]
即 \(\left(2+\frac{p}{2}\right)^2=16\). 因为 \(p>0\)，所以 \(2+\frac{p}{2}>0\)，从而 \(2+\frac{p}{2}=4\)，解得 \(p=4\).
\end{solution}

\begin{problem}
正六边形的中心和顶点共 \(7\text{ 个}\)点，以其中 \(3\text{ 个}\)点为顶点的三角形共有\fillinblank{}个. （用数字作答）
\end{problem}

\begin{answer}
\(32\)
\end{answer}

\begin{solution}
从中心和六个顶点中任取三个点共有 \(\mathrm C_7^3=35\) 种取法. 六个顶点中任意三个不共线；唯一的退化取法是中心与一对相对顶点共线，共有 \(3\) 条对角线，因此三角形共有
\[
\mathrm C_7^3-3=35-3=32
\]
个.
\end{solution}

\begin{problem}
\(\tan 20^{\circ} + \tan 40^{\circ} + \sqrt{3}\tan 20^{\circ}\tan 40^{\circ}\) 的值是\fillinblank{}.
\end{problem}

\begin{answer}
\(\sqrt{3}\)
\end{answer}

\begin{solution}
令 \(u=\tan20^{\circ}\)，\(v=\tan40^{\circ}\). 由正切加法公式
\[
\tan60^{\circ}=\frac{u+v}{1-uv}=\sqrt{3}
\]
所以 \(u+v+\sqrt{3}uv=\sqrt{3}\). 代回 \(u\)，\(v\) 得原式的值为 \(\sqrt{3}\).
\end{solution}

\begin{problem}
如图，正方形 \(ABCD\) 所在平面与正方形 \(ABEF\) 所在平面成 \(60^{\circ}\) 的二面角，则异面直线 \(AD\) 与 \(BF\) 所成角的余弦值是\fillinblank{}.

\bitmapfigure[max width=6.0cm,trim=0 0 0 0.45cm,clip]{img/1996/national_liberal/q19_fig1.png}
\end{problem}

\begin{answer}
\(\frac{\sqrt{2}}{4}\)
\end{answer}

\begin{solution}
设正方形边长为 \(a\)，取 \(AB\) 为公共棱. 令平面 \(ABCD\) 为水平面，并取
\(A=(0,0,0)\)，\(B=(a,0,0)\)，\(D=(0,a,0)\).
将平面 \(ABEF\) 绕 \(AB\) 与平面 \(ABCD\) 成 \(60^{\circ}\) 的二面角，则可取
\[
F=\left(0,\frac{a}{2},\frac{\sqrt{3}a}{2}\right)
\]
从而 \(\overrightarrow{AD}=(0,a,0)\)，\(\overrightarrow{BF}=\left(-a,\frac{a}{2},\frac{\sqrt{3}a}{2}\right)\). 因此
\[
\cos\angle(AD,BF)=\frac{\overrightarrow{AD}\cdot\overrightarrow{BF}}{|AD|\,|BF|}
=\frac{a^2/2}{a\cdot a\sqrt{2}}=\frac{\sqrt{2}}{4}
\]
故填 \(\frac{\sqrt{2}}{4}\).
\end{solution}

\section{解答题}

\begin{problem}
解不等式：\(\log_a(x + 1 - a) > 1\).
\end{problem}

\begin{solution}
由对数的定义，必须有 \(a>0\) 且 \(a\ne1\)，并且真数满足 \(x+1-a>0\). 又因为 \(1=\log_a a\)，所以分两种情况讨论. 当 \(a>1\) 时，对数函数 \(\log_a x\) 在 \((0,+\infty)\) 上单调递增，故 \(x+1-a>a\)，解得 \(x>2a-1\)，此时解集为 \(\left(2a-1,+\infty\right)\).

当 \(0<a<1\) 时，对数函数 \(\log_a x\) 在 \((0,+\infty)\) 上单调递减，故 \(0<x+1-a<a\)，解得 \(a-1<x<2a-1\)，此时解集为 \(\left(a-1,2a-1\right)\).
\end{solution}

\begin{problem}
设等比数列 \(\{a_{n}\}\) 的前 \(n\) 项和为 \(S_{n}\). 若 \(S_{3} + S_{6} = 2S_{9}\)，求数列的公比 \(q\).
\end{problem}

\begin{solution}
按等比数列的定义，\(a_{1}\ne0\) 且公比 \(q\ne0\). 若 \(q=1\)，则 \(S_{3}=3a_{1}\)，\(S_{6}=6a_{1}\)，\(S_{9}=9a_{1}\)，代入已知等式得 \(9a_{1}=18a_{1}\)，这与 \(a_{1}\ne0\) 矛盾. 因此 \(q\ne1\)，从而
\[
S_n=\frac{a_1\left(1-q^n\right)}{1-q}
\]
代入题设并约去非零因子 \(a_1/(1-q)\)，得 \(1-q^3+1-q^6=2\left(1-q^9\right)\)，即 \(q^3\left(2q^6-q^3-1\right)=0\). 因 \(q\ne0\)，且 \(2q^6-q^3-1=\left(2q^3+1\right)\left(q^3-1\right)\)，又 \(q\ne1\)，所以 \(q^3=-\frac{1}{2}\)，故
\[
q=-\sqrt[3]{\frac{1}{2}}=-\frac{\sqrt[3]{2}}{2}
\]
\end{solution}

\begin{problem}
已知 \(\triangle ABC\) 的三个内角 \(A\)，\(B\)，\(C\) 满足：\(A + C = 2B\)，\(\frac{1}{\cos A} + \frac{1}{\cos C} = -\frac{\sqrt{2}}{\cos B}\)，求 \(\cos \frac{A - C}{2}\) 的值.
\end{problem}

\begin{solution}
在三角形 \(ABC\) 中，\(A+B+C=\pi\). 由 \(A+C=2B\) 得 \(3B=\pi\)，所以 \(B=\frac{\pi}{3}\)，且 \(\frac{A+C}{2}=\frac{\pi}{3}\). 令 \(u=\frac{A+C}{2}\)，\(v=\frac{A-C}{2}\)，则 \(u=\frac{\pi}{3}\)，且 \(|v|<\frac{\pi}{3}\)，从而 \(\cos v>0\).

由和差化积公式，\(\cos A+\cos C=2\cos u\cos v=\cos v\)，并且 \(\cos A\cos C=\frac{\cos(A+C)+\cos(A-C)}{2}=\frac{\cos 2v-\frac{1}{2}}{2}\). 题设中的倒数有意义，所以 \(\cos A\ne0\)，\(\cos C\ne0\)，从而
\[
\frac{1}{\cos A}+\frac{1}{\cos C}=\frac{\cos v}{\left(\cos 2v-\frac{1}{2}\right)/2}=\frac{4\cos v}{4\cos^2v-3}
\]
又因 \(\cos B=\frac{1}{2}\)，题设给出上式等于 \(-2\sqrt{2}\). 令 \(c=\cos v\)，则 \(c>0\)，且
\[
\frac{4c}{4c^2-3}=-2\sqrt{2}\Longrightarrow 4\sqrt{2}c^2+2c-3\sqrt{2}=0
\]
解得 \(c=\frac{\sqrt{2}}{2}\) 或 \(c=-\frac{3\sqrt{2}}{4}\). 由于 \(c>0\)，所以
\[
\cos\frac{A-C}{2}=\cos v=\frac{\sqrt{2}}{2}
\]
\end{solution}

\begin{problem}
如图，在正三棱柱 \(ABC - A_{1}B_{1}C_{1}\) 中，\(AB = \frac{AA_{1}}{3} = a\)，\(E\)，\(F\) 分别是 \(BB_{1}\)，\(CC_{1}\) 上的点，\(BE = a\)，\(CF = 2a\).

\begin{enumerate}
\item 求证：面 \(AEF\) \(\perp\) 面 \(ACF\)；
\item 求三棱锥 \(A_{1} - AEF\) 的体积.
\end{enumerate}

\bitmapfigure[max width=6.0cm,trim=0 0.55cm 0 0,clip]{img/1996/national_liberal/q23_fig1.png}
\end{problem}

\begin{solution}
\begin{enumerate}
\item 建立空间直角坐标系，使底面 \(ABC\) 在 \(xOy\) 平面内，并取 \(A\left(0,0,0\right)\)，\(B\left(a,0,0\right)\)，\(C\left(\frac{a}{2},\frac{\sqrt{3}a}{2},0\right)\). 由 \(AA_1=3a\)、\(BE=a\)、\(CF=2a\)，得 \(E\left(a,0,a\right)\)，\(F\left(\frac{a}{2},\frac{\sqrt{3}a}{2},2a\right)\).

平面 \(AEF\) 的一个法向量为 \(\bs{n}_1=\overrightarrow{AE}\times\overrightarrow{AF}=\frac{a^2}{2}\left(-\sqrt{3},-3,\sqrt{3}\right)\)，平面 \(ACF\) 的一个法向量为 \(\bs{n}_2=\overrightarrow{AC}\times\overrightarrow{AF}=a^2\left(\sqrt{3},-1,0\right)\). 于是 \(\bs{n}_1\cdot\bs{n}_2=\frac{a^4}{2}\left(-3+3+0\right)=0\)，所以平面 \(AEF\) 与平面 \(ACF\) 垂直.

\item 由三棱锥体积的混合积公式，且 \(\overrightarrow{AA_1}=\left(0,0,3a\right)\)，有
\[
V_{A_1-AEF}=\frac{1}{6}\left|\left(\overrightarrow{AE}\times\overrightarrow{AF}\right)\cdot\overrightarrow{AA_1}\right|=\frac{1}{6}\left|\frac{a^2}{2}\left(-\sqrt{3},-3,\sqrt{3}\right)\cdot\left(0,0,3a\right)\right|=\frac{\sqrt{3}}{4}a^3
\]
\end{enumerate}
\end{solution}

\begin{problem}
某地现有耕地 \(10000\text{ 公顷}\)，规划 \(10\text{ 年}\)后粮食单产比现在增加 \(22\%\)，人均粮食占有量比现在提高 \(10\%\). 如果人口年增长率为 \(1\%\)，那么耕地平均每年至多只能减少多少公顷（精确到 \(1\text{ 公顷}\)）？（粮食单产 \(= \frac{\text{总产量}}{\text{耕地面积}}\)，人均粮食占有量 \(= \frac{\text{总产量}}{\text{总人口数}}\)）.
\end{problem}

\begin{solution}
设现在的粮食单产为 \(r\)，总人口数为 \(N\)，10 年后耕地面积为 \(S\)，平均每年减少耕地 \(x\) 公顷，则 \(S=10000-10x\). 10 年后单产为 \(1.22r\)，人口数为 \(1.01^{10}N\). 要使人均粮食占有量至少比现在提高 \(10\%\)，应有
\[
\frac{1.22rS}{1.01^{10}N}\geqslant1.10\frac{10000r}{N}
\]
由于 \(r>0\)、\(N>0\)，所以
\[
S\geqslant\frac{1.10\times10000\times1.01^{10}}{1.22}
\]
计算 \(1.01^2=1.0201\)，\(1.01^4=1.04060401\)，\(1.01^5=1.0510100501\)，从而 \(1.01^{10}=1.1046221254112045\ldots\). 因此
\(S\geqslant\frac{11000}{1.22}\times1.1046221254112045\ldots=9959.707688\ldots\)，\(x=\frac{10000-S}{10}\leqslant\frac{10000-9959.707688\ldots}{10}=4.029231\ldots\).
精确到 \(1\) 公顷，平均每年至多减少 \(4\) 公顷.
\end{solution}

\begin{problem}
已知 \(l_{1}\)，\(l_{2}\) 是过点 \(P\left(-\sqrt{2}, 0\right)\) 的两条互相垂直的直线，且 \(l_{1}\)，\(l_{2}\) 与双曲线 \(y^{2}-x^{2}=1\) 各有两个交点，分别为 \(A_{1}\)，\(B_{1}\) 和 \(A_{2}\)，\(B_{2}\).

\begin{enumerate}
\item 求 \(l_{1}\) 的斜率 \(k_{1}\) 的取值范围；
\item 若 \(A_{1}\) 恰是双曲线的一个顶点，求 \(\left|A_{2}B_{2}\right|\) 的值.
\end{enumerate}
\end{problem}

\begin{solution}
\begin{enumerate}
\item 先考虑过点 \(P\left(-\sqrt{2},0\right)\) 且斜率为 \(k\) 的直线 \(l\)：\(y=k\left(x+\sqrt{2}\right)\). 将它代入双曲线方程，得 \(\left(k^2-1\right)x^2+2\sqrt{2}k^2x+2k^2-1=0\). 当 \(k^2=1\) 时，方程退化为一次方程 \(2\sqrt{2}k^2x+2k^2-1=0\)，直线只有一个交点，因此必须有 \(k^2\ne1\). 此时方程的判别式为
\[
\Delta=\left(2\sqrt{2}k^2\right)^2-4\left(k^2-1\right)\left(2k^2-1\right)=4\left(3k^2-1\right)
\]
直线与双曲线有两个交点的充要条件是 \(\Delta>0\) 且 \(k^2\ne1\)，即 \(|k|>\frac{1}{\sqrt{3}}\) 且 \(k\ne\pm1\). 若直线为竖直线 \(x=-\sqrt{2}\)，它与双曲线的交点满足 \(y^2=3\)，确有两个交点；但其垂线为水平线 \(y=0\)，代入双曲线得 \(-x^2=1\)，无实交点. 因此在 \(l_1\)、\(l_2\) 均与双曲线有两个交点的条件下，二者均不是竖直线或水平线，均可用有限的非零斜率表示.

设 \(l_1\) 的斜率为 \(k_1\)，则 \(l_2\) 的斜率为 \(-\frac{1}{k_1}\). 对 \(l_2\) 应用同一条件，得 \(\left|\frac{1}{k_1}\right|>\frac{1}{\sqrt{3}}\)，即 \(|k_1|<\sqrt{3}\)，且仍有 \(k_1\ne\pm1\). 综合可得 \(l_1\) 的斜率范围为
\(\frac{1}{\sqrt{3}}<|k_1|<\sqrt{3}\)，\(k_1\ne\pm1\).

\item 双曲线 \(y^2-x^2=1\) 的两个顶点为 \(\left(0,1\right)\)、\(\left(0,-1\right)\). 若 \(A_1\) 是顶点，则由 \(P\left(-\sqrt{2},0\right)\) 到任一顶点的斜率可知 \(k_1=\pm\frac{1}{\sqrt{2}}\)，故 \(l_2\) 的斜率 \(m=-\frac{1}{k_1}\) 满足 \(m^2=2\). 将 \(y=m\left(x+\sqrt{2}\right)\) 代入双曲线方程，得
\[
2\left(x+\sqrt{2}\right)^2-x^2=1\Longleftrightarrow x^2+4\sqrt{2}x+3=0
\]
两交点的横坐标之差的绝对值为 \(2\sqrt{5}\)，而同一直线上纵坐标之差的绝对值为 \(|m|\) 倍的横坐标之差，所以
\[
\left|A_2B_2\right|=2\sqrt{5}\sqrt{1+m^2}=2\sqrt{5}\sqrt{3}=2\sqrt{15}
\]
\end{enumerate}
\end{solution}
